\scanpage{89}
% archive foliation: 88
Signalons pour finir ces généralités, que la formation des duals
$D(G)$ commute à l'extension de la base, qui transforme donc groupes
réflexifs en groupes réflexifs.

Nous nous intéressons par la suite au cas où
$\uline C=(\mathrm{Sch})_{/S}$, catégorie
des préschémas sur $S$, et $I=\uline G_{m\,S}$,
le ``groupe multiplicatif sur $S$'',
cf exposé I. Pour tout groupe \deleted{commutatif} ordinaire $M$, nous considérerons
le $S$-groupe $M_S$. On voit aussitôt que pour tout schéma en groupes
$I$ sur $S$, on a un isomorphisme canonique
(fonctoriel en $M,I$ et compatible avec l'extension de la base):
\[
\operatorname{Hom}_{S\text{-gr}}(M_S,I)
=\operatorname{Hom}_{\mathrm{gr}}(M,I(I/S))
\]
Appliquant ceci à un $S'$ variable sur $S$, on trouve un isomorphisme
\inserted{fonctoriel}
\[
D(M_S)(S')\qquad\operatorname{Hom}_{\mathrm{gr}}(M,\uline G_m(S'))
\]
On retrouve le foncteur déjà considéré dans l'exposé I,
\inserted{aussi noté $D(M)$,}
qui est représentable \deleted{\unsure{par un}}
pour $M$ commutatif comme
\[
D(M)=D(M_S)=\operatorname{Spec}\uline O_S\,M
\]
où $\uline O_S\,M$ désigne l'Algèbre du groupe $M$ à coefficients dans
$\uline O_S$. (No\deleted{\unsure{On vit}}tons
\inserted{d'ailleurs dans le cas général}
que $D(M)$ ne change pas si on remplace $M$ par $M$ rendu
\deleted{\unsure{Définition 1.1. Un}}
abélien, de sorte \deleted{\unsure{qu'il n'y a pratiquement}}
qu'on ne perd rien en supposant $M$ commutatif).

\uline{Définition} 1.1. Un schéma en groupes \inserted{$G$} sur $S$ est dit diagonalisable
s'il est isomorphe à un schéma de la forme \deleted{\unsure{et}}
$D(M)=D(M_S)=\uline{\operatorname{Hom}}_{S\text{-gr}}(M_S,\uline G_m$
pour un groupe commutatif $M$ convenable. On dit que $G$ est localement
diagonalisable, si tout point de $S$ admet un voisinage ouvert $U$, tel que
$G\,U$ soit diagonalisable.

\deleted{\uline{Définition} 1.2. Un schéma en groupes $I$ sur $S$ est appelé un
\uline{réseau constant} (resp. un réseau \uline{localement constant}) s'il est isomorphe à
un schéma en}
